\heading{数学物理方法（下）-25春-期末}
\noteinfo{题目来源于赛艇，答案由AI生成。}

\begin{question}[伴算符与斜自伴]

给定算符
\[
\hat L=-\frac{d}{dx},
\qquad
D(\hat L)=\{u(x):0<x<\pi,\ u\in L^2[0,\pi],\ u(\pi)=\gamma u(0)\}.
\]
求 $\hat L$ 的伴算符 $\hat L^\dagger$;并问复数 $\gamma$ 满足什么条件时有 $\hat L^\dagger=-\hat L$，并求此条件下 $\hat L^\dagger$ 的所有本征值和相应本征函数。

\begin{solution}

对任意 $u\in D(\hat L)$、$v\in D(\hat L^\dagger)$，有
\[
\langle \hat Lu,v\rangle
=-\int_0^{\pi}u'(x)\overline{v(x)}\,dx
=\int_0^{\pi}u(x)\overline{v'(x)}\,dx-u(\pi)\overline{v(\pi)}+u(0)\overline{v(0)}.
\]
由 $u(\pi)=\gamma u(0)$ 可得边界项为
\[
 u(0)\bigl(\overline{v(0)}-\gamma\overline{v(\pi)}\bigr).
\]
要对任意 $u$ 消失，必须有
\[
 v(0)=\overline{\gamma}\,v(\pi).
\]
于是
\ansmath{\hat L^\dagger=\frac{d}{dx}},
\qquad
\boxed{D(\hat L^\dagger)=\{v:0<x<\pi,\ v\in L^2[0,\pi],\ v(0)=\overline\gamma\,v(\pi)\}.}

因为 $-\hat L=d/dx$，要使 $\hat L^\dagger=-\hat L$，只需两个定义域一致，即
\[
 v(\pi)=\gamma v(0)
\quad\Longleftrightarrow\quad
 v(0)=\overline\gamma\,v(\pi).
\]
这当且仅当
\ansmath{|\gamma|=1.}

令 $\gamma=e^{i\theta}$（$\theta\in\mathbb R$），求 $\hat L^\dagger$ 的本征值：
\[
 v'=\lambda v
 \Longrightarrow v(x)=Ce^{\lambda x}.
\]
代入边界条件 $v(0)=\overline\gamma v(\pi)=e^{-i\theta}v(\pi)$ 得
\[
1=e^{-i\theta}e^{\lambda\pi}
\Longrightarrow
 e^{\lambda\pi}=e^{i\theta}.
\]
故
\ansmath{\lambda_n=\frac{i(\theta+2\pi n)}{\pi},\qquad n\in\mathbb Z,}
对应本征函数可取
\ansmath{v_n(x)=e^{\lambda_n x}.}

\end{solution}
\end{question}

\begin{question}[一维无限长杆上的热传导]

求解
\[
\begin{cases}
\dfrac{\partial u}{\partial t}-\kappa\dfrac{\partial^2u}{\partial x^2}=e^{-\alpha t},
&-\infty<x<+\infty,\ t>0,\\
u|_{t=0}=\delta(x).
\end{cases}
\]

\begin{solution}

把解分成两部分：
\[
 u=u_1+u_2.
\]
其中 $u_1$ 为热核对应于初值 $\delta(x)$ 的解：
\[
 u_1(x,t)=\frac{1}{\sqrt{4\pi\kappa t}}\exp\!\left(-\frac{x^2}{4\kappa t}\right).
\]

对强迫项，由于右端仅依赖于 $t$，可令 $u_2=u_2(t)$，则
\[
 u_2'(t)=e^{-\alpha t},\qquad u_2(0)=0.
\]
故
\[
 u_2(t)=\frac{1-e^{-\alpha t}}{\alpha}.
\]
所以
\ansmath{u(x,t)=\frac{1}{\sqrt{4\pi\kappa t}}\exp\!\left(-\frac{x^2}{4\kappa t}\right)
 +\frac{1-e^{-\alpha t}}{\alpha}.}

\end{solution}
\end{question}

\begin{question}[二维半圆形薄膜振动的 Green 函数]

求半圆域
\[
 y>0,\qquad x^2+y^2<R^2
\]
中波动方程
\[
\begin{cases}
\dfrac{\partial^2u}{\partial t^2}-a^2\nabla^2u=f(x,y,t),\\
u|_{y=0}=g(x,y,t),\qquad u|_{x^2+y^2=R^2}=h(x,y,t),\\
u|_{t=0}=\phi(x,y),\qquad \left.\dfrac{\partial u}{\partial t}\right|_{t=0}=\psi(x,y)
\end{cases}
\]
的 Green 函数。提示：
\[
\delta(x-x')\delta(y-y')=\frac1\rho\delta(\rho-\rho')\delta(\varphi-\varphi').
\]

\begin{solution}

在极坐标下，半圆域对应于
\[
 0<\rho<R,\qquad 0<\varphi<\pi.
\]
若按齐次 Dirichlet 边界讨论 Green 函数，则角向本征函数为
\[
 \sin(n\varphi),\qquad n=1,2,\dots
\]
径向本征函数为
\[
 J_n\!\Bigl(\frac{j_{n,m}}{R}\rho\Bigr),
\qquad m=1,2,\dots,
\]
其中 $j_{n,m}$ 是 $J_n$ 的第 $m$ 个正零点。对应特征值
\[
 \lambda_{n,m}=\frac{j_{n,m}}{R}.
\]
归一化模方为
\[
N_{n,m}=
\int_0^{\pi}\!\sin^2(n\varphi)\,d\varphi
\int_0^R \rho J_n^2\!\Bigl(\frac{j_{n,m}}{R}\rho\Bigr)\,d\rho
=\frac{\pi R^2}{4}J_{n+1}(j_{n,m})^2.
\]
因此 Green 函数可写成
\ansmath{G(\rho,\varphi,t;\rho',\varphi',t')
=H(t-t')
\sum_{n=1}^{\infty}\sum_{m=1}^{\infty}
\frac{\sin\bigl(a\lambda_{n,m}(t-t')\bigr)}{a\lambda_{n,m}N_{n,m}}
J_n\!\Bigl(\frac{j_{n,m}}{R}\rho\Bigr)
J_n\!\Bigl(\frac{j_{n,m}}{R}\rho'\Bigr)
\sin(n\varphi)\sin(n\varphi').}
将 $N_{n,m}$ 代入，可进一步写成
\[
G=H(t-t')\sum_{n,m}
\frac{4}{\pi a R j_{n,m}J_{n+1}(j_{n,m})^2}
\sin\Bigl(\frac{a j_{n,m}}{R}(t-t')\Bigr)
J_n\!\Bigl(\frac{j_{n,m}}{R}\rho\Bigr)
J_n\!\Bigl(\frac{j_{n,m}}{R}\rho'\Bigr)
\sin(n\varphi)\sin(n\varphi').
\]

\end{solution}
\end{question}

\begin{question}[泛函极值曲线]

求泛函
\[
J[u]=\iiint_{z>0,\ x^2+y^2+z^2<a^2}\bigl((\nabla u)^2-k^2u^2-zu\bigr)\,dxdydz
\]
的极值曲线 $u(x,y,z)$，其中
\[
 u|_{z=0}=0,
 \qquad u|_{x^2+y^2+z^2=a^2}=0,
\]
且 $a,k>0$ 并满足 $ka\ne \tan ka$。

\begin{solution}

Euler--Lagrange 方程为
\[
\frac{\partial L}{\partial u}-\nabla\cdot\frac{\partial L}{\partial (\nabla u)}=0
\quad\Longrightarrow\quad
\Delta u+k^2u=-\frac{z}{2}.
\]
右端与区域都关于 $z$ 轴对称，且源项正比于 $z=r\cos\theta$，因此设
\[
 u(r,\theta)=v(r)\cos\theta.
\]
一个显然的特解是
\[
 u_p=-\frac{z}{2k^2}.
\]
齐次方程
\[
(\Delta+k^2)u_h=0
\]
中，与 $\cos\theta$ 相应的正则解为球 Bessel 函数
\[
 u_h=A j_1(kr)\cos\theta.
\]
故总解为
\[
 u=\Bigl(Aj_1(kr)-\frac{r}{2k^2}\Bigr)\cos\theta.
\]
由球面边界 $u(a,\theta)=0$ 得
\[
 A=\frac{a}{2k^2j_1(ka)}.
\]
于是
\ansmath{u(r,\theta)=\left[\frac{a}{2k^2j_1(ka)}j_1(kr)-\frac{r}{2k^2}\right]\cos\theta.}
因 $z=r\cos\theta$，也可写成
\ansmath{u(x,y,z)=\frac{z}{2k^2}\left[\frac{a}{r}\frac{j_1(kr)}{j_1(ka)}-1\right],
 \qquad r=\sqrt{x^2+y^2+z^2}.}
其中
\[
 j_1(x)=\frac{\sin x-x\cos x}{x^2}.
\]
条件 $ka\ne \tan ka$ 恰好保证 $j_1(ka)\ne 0$。

\end{solution}
\end{question}

\clearpage

